\documentclass[book.tex]{subfiles}

\begin{document}
\chapter{Wick's Theorem}

\label{chap:wicks_theorem}
\noindent\rule{11cm}{0.4pt} \\
\textbf{Prerequisites:} \autoref{chap:qft_basics}  \\
\textbf{Difficulty Level:} ** \\
\noindent\rule{11cm}{0.4pt}
\vspace{5mm}

%\textcolor{red}{TODO: Should this be moved to the mathsection. Cite Shankar}

%Suppose that we have $N$ real random variables that form a vector 

%\begin{align}
%    |X\rangle &= [x_1,\dotsc,x_n],\quad x_i \in \mathbb{R}
%\end{align}
%We define a partition function over these variables by
%\begin{align}
%    Z(J) &= \prod_{i=1}^N \int_{-\infty}^{\infty} dx_i e^{-S_0(X) + \langle J | X \rangle}
%\end{align}
%Here $J$ is a source field vector
%\begin{align}
%    |J\rangle &= [J_1, \dotsc, J_n] \\
%    \langle J | X \rangle &= \sum_{i=1}^N J_i x_i
%\end{align}
%Recall that for a $1$-dimensional example, we %have the Gaussian integral
%\begin{align}
%    \int_{-\infty}^{\infty} e^{-\frac{1}{2}mx^2 + Jx} &= \sqrt{\frac{2\pi}{m}} \exp \Bigg [ %\frac{J^2}{2m} \Bigg ]
%\end{align}
%Let's now consider the matrix action
%\begin{align}
%    S(J) &= -\frac{1}{2}\langle X | M | X \rangle %+ \langle J | X \rangle
%\end{align}
%Here let $M$ be a real symmetric matrix that has %orthonormal eigen vectors
%\begin{align}
%    M|\alpha \rangle &= m_{\alpha}|\alpha \rangle %\\
%    \langle \alpha | \beta \rangle &= \delta_{ab}
%\end{align}
%We can expand vectors $|X\rangle$ and $|J\rangle$ %in terms of this orthonormal basis
%\begin{align}
%    |X\rangle &= \sum_{\alpha}x_\alpha |\alpha %\rangle \\
%    |J\rangle &= \sum_{\alpha}J_\alpha |\alpha %\rangle \\
%\end{align}
%These two equations mean that
%\begin{align}
%    \langle J | X \rangle &= %\sum_{\alpha}J_\alpha x_\alpha
%%\end{align}
%
%Since the mapping $x_i \mapsto x_\alpha$ is orthogonal, the jacobian $\mathcal{J}$ is identity. We can rewrite the action $S$ now as
%\begin{align}
%    S(J) &= -\frac{1}{2}\langle X | M | X \rangle + \langle J | X \rangle \\
%    &= \sum_{\alpha} \left ( -\frac{1}{2}x_\alpha^2m_\alpha + J_\alpha x_{\alpha} \right )
%\end{align}
%We can now factor the partition function into a producct of Gaussian integrals
%\begin{align}
%    Z(J) &= \prod_{i=1}^N \int_{-\infty}^{\infty} e^{-S(J)} \\
%    &= \prod_{\alpha} \int_{-\infty}^{\infty} dx_\alpha \exp \Bigg [ \sum_{\alpha} \left ( -\frac{1}{2}m_\alpha x_\alpha^2 + J_\alpha x_\alpha \right ) \Bigg ] \\
%    &= \prod_{\alpha} \left ( \int_{-\infty}^{\infty} dx_{\alpha} \exp \Bigg [ -\frac{1}{2}m_\alpha x_\alpha^2 + J_\alpha x_\alpha \Bigg ] \right ) \\
%    &= \prod_{\alpha} \sqrt{\frac{2\pi}{m_\alpha}} e^{\frac{\frac{1}{2}J_\alpha\frac{1}{m_\alpha}J_\alpha}{2m_{\alpha}}} \\
%    &= \frac{(2\pi)^{N/2}}{\sqrt{\det M}} e^{\frac{1}{2}\langle J | M^{-1}| J \rangle}
%\end{align}

%\textcolor{red}{TODO: I need to introduce $W(J)$ before I can take the n-dimensional Gaussian integral. Add this into the introduction to the Ising model chapter}
%\end{document}